\documentclass[11pt,a4paper]{article}

% Essential Packages
\usepackage[utf8]{inputenc}
\usepackage[margin=1in]{geometry}
\usepackage{amsmath,amsfonts,amssymb}
\usepackage{graphicx}
\usepackage{booktabs}
\usepackage{multirow}
\usepackage{xcolor}
\usepackage{hyperref}
\usepackage{fancyhdr}
\usepackage{enumitem}
\usepackage{titlesec}
\usepackage{caption}

% Color Definitions
\definecolor{primaryblue}{RGB}{26, 86, 160}
\definecolor{darkslate}{RGB}{30, 41, 59}
\definecolor{linkblue}{RGB}{2, 132, 199}
\definecolor{lightgrey}{RGB}{245, 247, 250}

% Hyperref Setup
\hypersetup{
    colorlinks=true,
    linkcolor=primaryblue,
    citecolor=primaryblue,
    urlcolor=linkblue,
    pdfauthor={Shafikul Islam Marwan},
    pdftitle={Lab Final Report: CNN Image Classification with PyTorch - 220121}
}

% Header and Footer
\pagestyle{fancy}
\fancyhf{}
\fancyhead[L]{\small \textcolor{darkslate}{\textbf{CSE 3-2: AI \& ML Lab Final Report}}}
\fancyhead[R]{\small \textcolor{darkslate}{\textbf{Student ID: 220121}}}
\fancyfoot[C]{\thepage}
\renewcommand{\headrulewidth}{0.4pt}
\renewcommand{\footrulewidth}{0.4pt}

% Section Title Formatting
\titleformat{\section}
  {\normalfont\Large\bfseries\color{primaryblue}}{\thesection}{1em}{}
\titleformat{\subsection}
  {\normalfont\large\bfseries\color{darkslate}}{\thesubsection}{1em}{}
\titleformat{\subsubsection}
  {\normalfont\normalsize\bfseries\color{darkslate}}{\thesubsubsection}{1em}{}

% Document Body
\begin{document}

% -------------------------------------------------------------
% Title and Student Metadata Block
% -------------------------------------------------------------
\begin{center}
    {\LARGE \textbf{\color{primaryblue} Convolutional Neural Network for Image Classification \\ with PyTorch \& Real-World Smartphone Validation}} \\[0.8em]
    {\large \textbf{Course:} Artificial Intelligence \& Machine Learning Lab (Lab Final)} \\[0.5em]
    \rule{\linewidth}{1pt} \\[0.8em]
    
    \begin{tabular}{ll}
        \textbf{Student Name:} & Shafikul Islam Marwan \\
        \textbf{Student ID:} & 220121 \\
        \textbf{Department:} & Computer Science and Engineering \\
        \textbf{Dataset Chosen:} & Hand Gestures (Rock-Paper-Scissors Dataset) \\
        \textbf{Framework:} & PyTorch 2.0+ \& Torchvision \\
        \textbf{Execution Environment:} & Google Colab (Automated Pipeline with GPU/CPU support) \\
        \textbf{GitHub Repository:} & \url{https://github.com/Marwanthe0/CSE/tree/main/Third%20Year/3-2/Artificial%20Intelligence%20and%20ML/Lab%20Final} \\
        \textbf{Google Colab Link:} & \url{https://colab.research.google.com/drive/1lCTyW3xMkFSMuc5Su9Zgi0NG4J7OAbLm?usp=sharing} \\
    \end{tabular}
    \\[0.8em]
    \rule{\linewidth}{1pt}
\end{center}

% -------------------------------------------------------------
% Abstract
% -------------------------------------------------------------
\begin{abstract}
\noindent In this laboratory assignment, an end-to-end deep learning pipeline for image classification is developed using the PyTorch framework. A custom Convolutional Neural Network (CNN) is designed from scratch, trained, and evaluated on Laurence Moroney's standard Rock-Paper-Scissors (RPS) dataset containing 2,520 training and 372 test images. To bridge standard benchmark training with practical application, 10 real-world smartphone photographs corresponding to the dataset classes (Rock, Paper, Scissors) were captured and processed through the exact same transformation pipeline. The system is designed to be fully automated on Google Colab, autonomously retrieving custom assets from a public GitHub repository and downloading benchmark datasets without requiring manual runtime intervention. On the standard benchmark test set, the trained CNN achieved an overall classification accuracy of \textbf{95.43\%} with an F1-score of 0.95. Real-world smartphone testing was conducted, and the domain shift between synthetic benchmark data and real-world camera captures is analyzed in detail.
\end{abstract}

% -------------------------------------------------------------
% 1. Introduction & Objectives
% -------------------------------------------------------------
\section{Introduction and Objectives}
Image classification represents a foundational problem in computer vision and artificial intelligence. While classical machine learning algorithms rely on handcrafted feature extractors (e.g., SIFT, HOG), Convolutional Neural Networks (CNNs) learn hierarchical representations directly from raw pixel data through alternating convolution, non-linear activation, and spatial downsampling operations.

The primary objectives of this laboratory assignment are:
\begin{enumerate}[leftmargin=1.5em, itemsep=2pt]
    \item \textbf{Custom CNN Architecture:} Design and implement a multi-block CNN by subclassing PyTorch's \texttt{nn.Module}, incorporating Convolution, Batch Normalization, ReLU activation, Max Pooling, and Dropout regularization.
    \item \textbf{End-to-End Pipeline:} Implement data ingestion, augmentation, normalization, and minibatch iteration using \texttt{torchvision.transforms} and \texttt{DataLoader}.
    \item \textbf{Training and Optimization:} Train the model from scratch using Cross-Entropy Loss and the Adam optimizer, logging training and validation performance metrics across all epochs.
    \item \textbf{Comprehensive Evaluation:} Evaluate the trained model via Training Curves (Loss vs. Epochs, Accuracy vs. Epochs), Confusion Matrix heatmap, Classification Report, and Visual Error Analysis.
    \item \textbf{Real-World Validation (The Phone Task):} Capture and classify 10 independent smartphone images to assess model generalization under natural lighting and real-world background variations.
    \item \textbf{Full Reproducibility:} Implement an automated pipeline capable of executing entirely via a single ``Run All'' trigger in Google Colab through seamless GitHub synchronization.
\end{enumerate}

% -------------------------------------------------------------
% 2. Dataset and Preprocessing Pipeline
% -------------------------------------------------------------
\section{Dataset and Preprocessing Pipeline}

\subsection{Standard Benchmark Dataset}
The primary dataset used for model training and evaluation is Laurence Moroney's standard \textbf{Rock-Paper-Scissors (RPS)} dataset. The dataset comprises synthetically rendered hand gestures with diverse skin tones, poses, and angles against clean, uniform white backgrounds.
\begin{itemize}[leftmargin=1.5em, itemsep=2pt]
    \item \textbf{Training Set:} 2,520 RGB images ($300 \times 300$ resolution), evenly split into 840 images per class: \textit{Paper}, \textit{Rock}, and \textit{Scissors}.
    \item \textbf{Test/Validation Set:} 372 RGB images, containing 124 images per class.
\end{itemize}

\subsection{Custom Smartphone Dataset (The Phone Task)}
In compliance with assignment requirements, 10 real-world photographs were captured using a personal smartphone camera under ambient indoor lighting:
\begin{itemize}[leftmargin=1.5em, itemsep=2pt]
    \item \textbf{Distribution:} 4 Rock images (\texttt{rock\_1.jpeg} to \texttt{rock\_4.jpeg}), 3 Paper images (\texttt{paper\_1.jpeg} to \texttt{paper\_3.jpeg}), and 3 Scissors images (\texttt{scissors\_1.jpeg} to \texttt{scissors\_3.jpeg}).
    \item \textbf{Storage:} Uploaded to the project's \texttt{dataset/} folder in the public GitHub repository.
\end{itemize}

\subsection{Data Preprocessing and Augmentation Pipeline}
To ensure computational efficiency and regularize the network against spatial variance, all images are dynamically transformed during batch loading:
\begin{itemize}[leftmargin=1.5em, itemsep=2pt]
    \item \textbf{Training Transforms:} 
    \begin{align*}
        T_{\text{train}} = \text{Resize}(150 \times 150) \circ \text{RandomHorizontalFlip}(p=0.5) \circ \text{RandomRotation}(10^\circ) \\
        \circ\, \text{ColorJitter}(\text{brightness}=0.1, \text{contrast}=0.1) \circ \text{ToTensor}() \circ \text{Normalize}(\mu, \sigma)
    \end{align*}
    where $\mu = [0.485, 0.456, 0.406]$ and $\sigma = [0.229, 0.224, 0.225]$ (standard ImageNet statistics).
    \item \textbf{Validation and Custom Image Transforms:}
    \begin{align*}
        T_{\text{eval}} = \text{Resize}(150 \times 150) \circ \text{ToTensor}() \circ \text{Normalize}(\mu, \sigma)
    \end{align*}
\end{itemize}
The exact same normalization and resizing dimensions are applied to the 10 custom phone images to ensure mathematical consistency.

% -------------------------------------------------------------
% 3. CNN Model Architecture
% -------------------------------------------------------------
\section{Model Architecture}
The classification model is implemented as a custom subclass of \texttt{nn.Module}. The network consists of three sequential convolutional feature extraction blocks followed by a fully connected classification head with adaptive pooling and dropout layers.

\begin{table}[h!]
\centering
\small
\caption{Layer-by-Layer Architectural Specifications of the Custom CNN}
\label{tab:architecture}
\begin{tabular}{@{}llcc@{}}
\toprule
\textbf{Block / Layer} & \textbf{Component Specifications} & \textbf{Output Shape} & \textbf{Parameters} \\ \midrule
\textbf{Input Layer} & RGB Normalized Image & $3 \times 150 \times 150$ & 0 \\
\textbf{Conv Block 1} & $\text{Conv2D}(3 \to 32, k=3, p=1) + \text{BatchNorm2D}(32) + \text{ReLU} + \text{MaxPool2D}(2,2)$ & $32 \times 75 \times 75$ & 960 \\
\textbf{Conv Block 2} & $\text{Conv2D}(32 \to 64, k=3, p=1) + \text{BatchNorm2D}(64) + \text{ReLU} + \text{MaxPool2D}(2,2)$ & $64 \times 37 \times 37$ & 18,624 \\
\textbf{Conv Block 3} & $\text{Conv2D}(64 \to 128, k=3, p=1) + \text{BatchNorm2D}(128) + \text{ReLU} + \text{MaxPool2D}(2,2)$ & $128 \times 18 \times 18$ & 74,112 \\
\textbf{Pooling} & $\text{AdaptiveAvgPool2D}((4, 4)) + \text{Flatten}()$ & $2048$ & 0 \\
\textbf{Dropout 1} & $\text{Dropout}(p=0.4)$ & $2048$ & 0 \\
\textbf{FC Layer 1} & $\text{Linear}(2048 \to 256) + \text{ReLU}() + \text{Dropout}(p=0.3)$ & $256$ & 524,544 \\
\textbf{Output Layer} & $\text{Linear}(256 \to 3)$ & $3$ & 771 \\ \midrule
\textbf{Total Parameters} & \multicolumn{3}{r}{\textbf{619,011 (All Trainable)}} \\ \bottomrule
\end{tabular}
\end{table}

\noindent \textbf{Mathematical Formulation of the Forward Pass:}
For an input tensor $X \in \mathbb{R}^{B \times 3 \times 150 \times 150}$, the feature map transformation for block $i$ is:
\begin{equation}
    H_i = \text{MaxPool}\Big(\text{ReLU}\big(\text{BatchNorm}(\text{Conv2d}(H_{i-1}))\big)\Big), \quad i \in \{1, 2, 3\}
\end{equation}
The spatial representations are pooled into a fixed dimension $z = \text{AdaptiveAvgPool2D}(H_3) \in \mathbb{R}^{B \times 128 \times 4 \times 4}$, flattened to a vector $\hat{z} \in \mathbb{R}^{B \times 2048}$, and mapped to class logits $\mathbf{z} \in \mathbb{R}^{B \times 3}$:
\begin{equation}
    \mathbf{z} = W_2 \cdot \text{Dropout}\big(\text{ReLU}(W_1 \hat{z} + b_1)\big) + b_2
\end{equation}
The class prediction probabilities are obtained using the Softmax function:
\begin{equation}
    P(Y = c \mid X) = \frac{e^{z_c}}{\sum_{j=1}^{3} e^{z_j}}, \quad c \in \{\text{Paper}, \text{Rock}, \text{Scissors}\}
\end{equation}

% -------------------------------------------------------------
% 4. Training Methodology and Hyperparameters
% -------------------------------------------------------------
\section{Training Methodology and Experimental Setup}
The model was trained from scratch using empirical cross-entropy loss over 10 epochs. The optimization objective is given by:
\begin{equation}
    \mathcal{L}_{\text{CE}}(\theta) = -\frac{1}{N}\sum_{i=1}^{N} \log P\big(Y = y_i \mid X_i; \theta\big)
\end{equation}
where $\theta$ denotes the trainable weights of the network.

\begin{table}[h!]
\centering
\small
\caption{Hyperparameter Configuration and Training Setup}
\label{tab:hyperparams}
\begin{tabular}{@{}ll|ll@{}}
\toprule
\textbf{Hyperparameter} & \textbf{Value} & \textbf{Parameter} & \textbf{Value} \\ \midrule
Batch Size ($B$) & 64 & Loss Function & CrossEntropyLoss \\
Initial Learning Rate ($\alpha$) & $1 \times 10^{-3}$ & Number of Epochs & 10 \\
Optimizer & Adam ($\beta_1=0.9, \beta_2=0.999$) & Weight Initialization & PyTorch Kaiming Default \\
Input Resolution & $150 \times 150 \times 3$ & Device & PyTorch Auto-detection (CPU/CUDA) \\
\bottomrule
\end{tabular}
\end{table}

% -------------------------------------------------------------
% 5. Results and Performance Analysis
% -------------------------------------------------------------
\section{Experimental Results and Performance Analysis}

\subsection{Training and Validation Progression}
Table~\ref{tab:epoch_history} details the empirical loss and accuracy progression across the 10 training epochs. The model converged rapidly, achieving $>95\%$ training accuracy by Epoch 2 and maintaining high validation accuracy throughout training.

\begin{table}[h!]
\centering
\small
\caption{Epoch-by-Epoch Loss and Accuracy Progression}
\label{tab:epoch_history}
\begin{tabular}{@{}ccccc@{}}
\toprule
\textbf{Epoch} & \textbf{Training Loss} & \textbf{Training Accuracy (\%)} & \textbf{Validation Loss} & \textbf{Validation Accuracy (\%)} \\ \midrule
1 & 0.6223 & 74.60\% & 0.3053 & 83.33\% \\
2 & 0.1468 & 95.87\% & 0.1956 & 93.82\% \\
3 & 0.0957 & 97.62\% & 1.1049 & 68.55\% \\
4 & 0.0833 & 97.34\% & 0.3243 & 91.40\% \\
5 & 0.0800 & 97.70\% & 0.4889 & 82.80\% \\
6 & 0.0360 & 99.01\% & 0.1203 & 94.35\% \\
7 & 0.0241 & 99.37\% & 1.0759 & 78.49\% \\
8 & 0.0200 & 99.37\% & 0.2145 & 93.01\% \\
9 & 0.0235 & 99.17\% & 0.2970 & 92.74\% \\
\textbf{10} & \textbf{0.0261} & \textbf{99.13\%} & \textbf{0.1396} & \textbf{95.43\%} \\
\bottomrule
\end{tabular}
\end{table}

\subsection{Evaluation on Standard Test Set}
The trained model was evaluated on the independent 372-image test set. The classification report is summarized in Table~\ref{tab:metrics}.

\begin{table}[h!]
\centering
\small
\caption{Per-Class Classification Metrics on Standard Test Set}
\label{tab:metrics}
\begin{tabular}{@{}lcccc@{}}
\toprule
\textbf{Class} & \textbf{Precision} & \textbf{Recall} & \textbf{F1-Score} & \textbf{Support} \\ \midrule
\textbf{Paper} & 1.00 & 0.86 & 0.93 & 124 \\
\textbf{Rock} & 0.95 & 1.00 & 0.97 & 124 \\
\textbf{Scissors} & 0.93 & 1.00 & 0.96 & 124 \\ \midrule
\textbf{Overall Accuracy} & \multicolumn{3}{c}{\textbf{95.43\%}} & \textbf{372} \\
\textbf{Macro Average} & 0.96 & 0.95 & 0.95 & 372 \\
\textbf{Weighted Average} & 0.96 & 0.95 & 0.95 & 372 \\ \bottomrule
\end{tabular}
\end{table}

\noindent \textbf{Confusion Matrix Summary:}
The model misclassified only 17 out of 372 test images ($4.57\%$ total error rate). All misclassifications were confined to the \textit{Paper} class being predicted as \textit{Rock} or \textit{Scissors}, whereas \textit{Rock} and \textit{Scissors} achieved a perfect $100\%$ recall on the test set.

% -------------------------------------------------------------
% 6. Real-World Smartphone Testing (The Phone Task)
% -------------------------------------------------------------
\section{Real-World Smartphone Testing and Domain Shift Analysis}

\subsection{Prediction Results on Custom Smartphone Images}
The 10 custom smartphone photographs stored in \texttt{dataset/} were processed through the forward inference pipeline. The model generated the predictions reported in Table~\ref{tab:custom_pred}.

\begin{table}[h!]
\centering
\small
\caption{Inference Results on Custom Smartphone Photographs}
\label{tab:custom_pred}
\begin{tabular}{@{}lccc@{}}
\toprule
\textbf{Filename} & \textbf{Ground Truth} & \textbf{Predicted Class} & \textbf{Model Confidence (\%)} \\ \midrule
\texttt{rock\_1.jpeg} & Rock & Rock & 100.0\% \\
\texttt{rock\_2.jpeg} & Rock & Rock & 100.0\% \\
\texttt{rock\_3.jpeg} & Rock & Rock & 100.0\% \\
\texttt{rock\_4.jpeg} & Rock & Rock & 100.0\% \\
\texttt{paper\_1.jpeg} & Paper & Rock & 99.9\% \\
\texttt{paper\_2.jpeg} & Paper & Rock & 100.0\% \\
\texttt{paper\_3.jpeg} & Paper & Rock & 97.6\% \\
\texttt{scissors\_1.jpeg} & Scissors & Rock & 100.0\% \\
\texttt{scissors\_2.jpeg} & Scissors & Rock & 100.0\% \\
\texttt{scissors\_3.jpeg} & Scissors & Rock & 100.0\% \\ \bottomrule
\end{tabular}
\end{table}

\subsection{Domain Shift and Error Analysis}
While the model demonstrates high performance on the standard test set ($95.43\%$), it exhibits a strong prediction bias toward the \textit{Rock} class when applied to real-world smartphone photos. This discrepancy is a classic demonstration of \textbf{Domain Shift} and \textbf{Out-of-Distribution (OOD) Covariate Shift}:
\begin{enumerate}[leftmargin=1.5em, itemsep=2pt]
    \item \textbf{Background Invariance Gap:} The training dataset features pure synthetic white backgrounds with sharp edges. Real-world camera images introduce subtle ambient illumination gradients, shadows, and surface texture.
    \item \textbf{Skin Tone and Hand Geometry Variance:} The synthetic training samples possess specific texture renderings, whereas smartphone captures introduce camera sensor post-processing, lens distortion, and varied flesh tones.
    \item \textbf{Mitigation Strategies:} Domain generalization can be enhanced in future work through extensive background augmentation (e.g., CutMix, Mosaic, Random Erasing), HSV color-space transformations, and fine-tuning on a small real-world support set.
\end{enumerate}

% -------------------------------------------------------------
% 7. Reproducibility and Automation
% -------------------------------------------------------------
\section{Reproducibility and Repository Setup}
The project strictly satisfies all automation and reproducibility constraints specified in the assignment:
\begin{itemize}[leftmargin=1.5em, itemsep=2pt]
    \item \textbf{Automated GitHub Synchronization:} The Google Colab notebook executes \texttt{!git clone} directly to download the 10 custom phone images and project structure without manual file uploading.
    \item \textbf{Self-Contained Data Loading:} Benchmark training and testing datasets are fetched dynamically via direct zip extraction and loaded using \texttt{torchvision.datasets.ImageFolder}.
    \item \textbf{State Dictionary Storage:} The trained model state dictionary is saved as \texttt{220121.pth} in the \texttt{model/} directory.
    \item \textbf{One-Click Execution:} Opening the Colab link and selecting \textbf{Runtime $\to$ Run All} executes the entire pipeline sequentially: cloning repository assets, downloading dataset, training the CNN, generating all plots/visualizations, and evaluating real-world smartphone samples.
\end{itemize}

% -------------------------------------------------------------
% 8. Conclusion
% -------------------------------------------------------------
\section{Conclusion}
In this assignment, a complete PyTorch deep learning workflow was successfully implemented for hand gesture classification. A custom 3-block CNN comprising 619,011 parameters was constructed, trained from scratch, and achieved a test accuracy of \textbf{95.43\%} on benchmark data. Real-world testing with smartphone images illustrated practical domain shift challenges in machine learning deployment. All code, datasets, trained model weights (\texttt{220121.pth}), executed notebooks, and documentation are publicly accessible on GitHub and Google Colab.

% -------------------------------------------------------------
% References
% -------------------------------------------------------------
\begin{thebibliography}{9}
\bibitem{pytorch}
Paszke, A., et al. (2019). ``PyTorch: An Imperative Style, High-Performance Deep Learning Library.'' \textit{Advances in Neural Information Processing Systems (NeurIPS)}, 32.
\bibitem{rps_dataset}
Moroney, L. (2019). ``Rock, Paper, Scissors Dataset.'' \textit{TensorFlow Datasets}. \url{https://storage.googleapis.com/download.tensorflow.org/data/rps.zip}
\bibitem{adam}
Kingma, D. P., \& Ba, J. (2014). ``Adam: A Method for Stochastic Optimization.'' \textit{arXiv preprint arXiv:1412.6980}.
\end{thebibliography}

\end{document}
